\documentclass[11pt,a4paper]{article}
\usepackage[margin=2.3cm]{geometry}
\usepackage{amsmath,amssymb}
\usepackage{booktabs}
\usepackage{graphicx}
\usepackage{microtype}
\usepackage[raggedright]{titlesec}
\usepackage{enumitem}
\usepackage{float}
\usepackage{caption}
\usepackage[table]{xcolor}
\usepackage{tikz}
\usepackage[hidelinks]{hyperref}
\usetikzlibrary{arrows.meta,positioning}
\graphicspath{{../paper/figures/}}
\setlist{itemsep=2pt,topsep=4pt}
\setlength{\parindent}{0pt}
\setlength{\parskip}{6pt}

\begin{document}

{\Large\textbf{Progress report: quantum channel attention for sign recognition}}\par\medskip

This is a summary of what I have done so far on the project: the reading,
the model, the experiments, the results and the paper draft.

\section*{1. What I set out to do}

I wanted to check if a small quantum circuit can work as a channel-attention
block inside a CNN for American Sign Language (ASL) letter recognition. The
plan was to put an 8-qubit circuit into DenseNet121 and see if it helps
accuracy or saves parameters compared with normal classical attention.

Short answer so far: it saves parameters compared with SE attention, but it
does not improve accuracy, and a classical layer of the same size does just as
well. So the result is negative.

\section*{2. Background reading}

I went through around 28 papers (2020--2025, all journal or conference
papers, no arXiv-only ones). Main points I took from them:

\begin{itemize}
  \item Deep CNNs already get close to 100\% on the Kaggle ASL alphabet
        dataset. Alsharif et al.\ (2023) got 99.98\% with ResNet-50. So that
        dataset alone cannot show whether a new module helps.
  \item The real problem in sign recognition is new signers, people the
        model never saw in training (Erdo\u{g}an and Bayta\c{s}, 2024). Sahoo
        et al.\ (2022) tested the MU HandImages dataset with
        leave-one-subject-out testing, so I used the same idea.
  \item Hybrid quantum-classical image models usually put a circuit after a
        classical feature extractor (Mari et al., 2020; Long et al., 2025).
        Budiutama et al.\ (2024) have ``channel attention'' for quantum CNNs,
        but it works on quantum states, not on CNN image channels.
  \item Theory papers say small circuits are limited in what they can
        represent (Schuld et al., 2021), classical models with data can often
        match them (Huang et al., 2021), and training can suffer from
        vanishing gradients (Cerezo et al., 2021).
  \item For attention: SE (Hu et al., 2020) is the usual baseline, and
        ECA-Net (Wang et al., 2020) gets similar results with fewer
        parameters.
  \item Several papers say quantum models need fair classical baselines
        (Raubitzek and Mallinger, 2023; Schuld and Killoran, 2022), and that
        random seeds can change which method looks best (Bouthillier et al.,
        2021). Because of this I added a classical layer with the same number
        of parameters and repeated runs with different seeds.
\end{itemize}

\section*{3. What I built}

The model is DenseNet121 with a quantum channel attention (QCA) gate between
the global pooling and the classifier:

\begin{enumerate}
  \item The 1,024 pooled features are projected down to 8 numbers and scaled
        with $\tanh$ to the range $(-\pi/2, \pi/2)$.
  \item These 8 numbers are loaded into 8 qubits with $R_Y$ rotations, passed
        through 2 entangling layers, and measured (Pauli-$Z$).
  \item The 8 outputs are projected back up to 1,024 values and passed
        through a sigmoid to get one weight per channel.
  \item The features are multiplied by these weights before the classifier.
\end{enumerate}

The circuit is simulated in PennyLane on CPU. The whole gate has 17,464
parameters (8,200 down-projection, 48 in the circuit, 9,216 up-projection).

\begin{figure}[H]
\centering
\resizebox{\textwidth}{!}{\input{../paper/figures/qca_pipeline.tex}}
\caption*{The model (a) and what is inside the QCA gate (b).}
\end{figure}

To compare fairly, I also implemented three other options in the same place:

\begin{table}[H]
\centering
\begin{tabular}{lrl}
\toprule
Module & Parameters & Why \\
\midrule
QCA (quantum) & 17,464 & the one being tested \\
Matched MLP & 17,480 & same shape as QCA, classical layer instead of circuit \\
SE attention & 131,072 & standard classical attention \\
No attention & 0 & plain DenseNet121 \\
\bottomrule
\end{tabular}
\end{table}

\section*{4. Data}

\begin{itemize}
  \item \textbf{Kaggle ASL Alphabet:} 87,000 images, 29 classes. Split
        80/10/10 with a fixed seed (8,700 test images).
  \item \textbf{MU HandImages:} 2,515 images in total; I used the 1,815 A--Z
        images from 5 signers. The signers have 650, 645, 130, 130 and 260
        images.
  \item No horizontal flipping, because flipping a hand changes the sign.
\end{itemize}

\begin{figure}[H]
\centering
\includegraphics[width=\textwidth]{dataset_samples.pdf}
\caption*{Some letters from Kaggle (top) and MU (bottom).}
\end{figure}

\section*{5. Experiments I ran, in order}

\begin{enumerate}
  \item \textbf{Baseline.} DenseNet121 fine-tuned on Kaggle: 100.00\% test
        accuracy. This showed the Kaggle split is too easy to compare models.
  \item \textbf{QCA on Kaggle, full training.} DenseNet121 + QCA and
        ResNet50 + QCA both got 100.00\% test accuracy. Same as baseline, so
        this says nothing about QCA.
  \item \textbf{Frozen backbone, full data (one seed).} Only the head trains,
        so the attention block matters more. SE 97.10\%, QCA 95.98\%, no
        attention 95.89\%.
  \item \textbf{Low-data runs.} Frozen backbone, 1\%, 5\% and 10\% of the
        training data, all four modules, three seeds each (36 runs).
  \item \textbf{Re-uploading test.} A QCA version that loads the input again
        before each circuit layer. One run at 1\% data: 61.38\% vs.\ 61.83\%
        for the normal QCA, so I did not continue with it.
  \item \textbf{MU unseen-signer test (main experiment).} As my supervisor
        asked, all layers are trained here. I train on 4 signers and test on
        the 5th, for each signer and each of the four modules (20 runs).
        15 epochs, batch 32, learning rate $10^{-5}$ for the backbone and
        $10^{-3}$ for the head, seed 42.
\end{enumerate}

\section*{6. Results}

\textbf{MU, unseen signers.} The main number is the average over the 5
signers (each signer counts equally).

\begin{table}[H]
\centering
\begin{tabular}{lcccc}
\toprule
Module & Avg.\ over signers (\%) & All images pooled (\%) & Macro-F1 (\%) & $p$ vs.\ QCA \\
\midrule
No attention & 94.78 $\pm$ 5.17 & 94.82 & 94.61 & 0.543 \\
Matched MLP & 94.55 $\pm$ 4.02 & 94.88 & 93.95 & 0.943 \\
QCA & 94.50 $\pm$ 4.95 & 94.99 & 94.26 & -- \\
SE & 93.72 $\pm$ 5.17 & 94.71 & 93.04 & 0.116 \\
\bottomrule
\end{tabular}
\end{table}

All four are within about 1 point and none of the differences is
significant (paired t-test over signers). QCA is slightly ahead on the pooled
number only because the bigger signers count more there.

Per signer:

\begin{table}[H]
\centering
\begin{tabular}{lccccc}
\toprule
Module & Signer 1 & Signer 2 & Signer 3 & Signer 4 & Signer 5 \\
\midrule
QCA & 95.38 & 97.52 & 96.92 & 96.92 & 85.77 \\
Matched MLP & 94.77 & 97.21 & 97.69 & 95.38 & 87.69 \\
SE & 95.69 & 97.52 & 95.38 & 95.38 & 84.62 \\
No attention & 95.38 & 96.59 & 98.46 & 97.69 & 85.77 \\
\bottomrule
\end{tabular}
\end{table}

Signer 5 is much harder for every model (about 85--88\%). The difference
between modules on any one signer is at most about 3 points, so who the
signer is matters much more than which attention block is used.

\begin{figure}[H]
\centering
\includegraphics[width=0.95\textwidth]{mu_efficiency.pdf}
\caption*{(a) MU accuracy vs.\ attention parameters. (b) QCA minus each other module, per signer.}
\end{figure}

Over all 1,815 MU images, QCA got 1,724 right. Most mistakes are between
similar hand shapes: M/N (16 and 13), V$\to$W (11), A$\to$T (6).

\begin{figure}[H]
\centering
\includegraphics[width=0.55\textwidth]{qca_confusion_matrix.pdf}
\caption*{QCA confusion matrix on MU, all five test signers together.}
\end{figure}

\begin{table}[H]
\raggedright\textbf{Kaggle, low data (frozen backbone, mean $\pm$ SD over 3 seeds):}\par\medskip
\centering
\begin{tabular}{lccc}
\toprule
Module & 1\% data & 5\% data & 10\% data \\
\midrule
SE & 66.17 $\pm$ 0.42 & 87.22 $\pm$ 0.88 & 91.12 $\pm$ 0.45 \\
No attention & 64.69 $\pm$ 0.63 & 86.43 $\pm$ 0.55 & 90.19 $\pm$ 0.60 \\
Matched MLP & 60.80 $\pm$ 1.31 & 86.39 $\pm$ 0.47 & 90.93 $\pm$ 0.24 \\
QCA & 62.16 $\pm$ 1.16 & 86.13 $\pm$ 0.93 & 90.36 $\pm$ 0.35 \\
\bottomrule
\end{tabular}
\end{table}

SE is best at every size. QCA beats no attention only at 10\%, and only by
0.17 points. QCA beats the matched MLP at 1\%, but the MLP is ahead again at
10\%.

\section*{7. What I conclude}

\begin{itemize}
  \item QCA uses about 7.5 times fewer attention parameters than SE (17,464
        vs.\ 131,072) and does not lose accuracy on new signers.
  \item But the quantum circuit itself does not seem to add anything. The
        classical MLP with the same size gives the same results, so the
        saving comes from the small 8-unit bottleneck, not from the quantum
        part.
  \item Limits: MU was run with one seed only, there are only 5 signers, and
        the circuit is simulated without noise. Differences under about
        1 point cannot be judged from this.
\end{itemize}

\section*{8. Paper so far}

I have an 8-page conference draft (target limit 10 pages) with 4 figures,
4 tables and 28 references. My supervisor has reviewed it four times, and I
made these changes:

\begin{itemize}
  \item Presented QCA-DenseNet as the one proposed method instead of a
        comparison of four, with three clear contributions that the results
        section checks one by one.
  \item Made all captions one line, put every figure and table after the
        text that discusses it, and kept one readable confusion matrix.
  \item Changed Limitations to Future Scope and removed everything after the
        Conclusion.
  \item Kept only real references from 2020 onwards, checked each one by DOI,
        removed the only arXiv-only paper, and numbered citations in order.
  \item Removed footnotes and personal words (I/we/our) from the paper.
  \item Cut the Abstract to about 150 words, the Conclusion to about 120,
        Future Scope to one sentence, and Discussion to two paragraphs.
  \item Fixed a wrong number: the MU dataset has 2,515 images, not 2,425.
\end{itemize}

I used Claude (an AI assistant) to help draft and edit the paper text. All
the numbers come from my experiment result files. Before submitting, I need
to rewrite the text in my own words or disclose the AI use, as the venue
requires.

\section*{9. What is left}

\begin{itemize}
  \item Discuss the negative result with my supervisor before running
        anything new.
  \item Decide if I should repeat the MU runs with more seeds.
  \item Check the licence of the MU dataset copy I used (the paper shows MU
        images).
  \item Pick the conference and its template, and finalise the title.
  \item Rewrite the text in my own words or disclose the AI help, then run a
        plagiarism check.
  \item Possible later work, if approved: compare with ECA-Net, test the
        circuit with noise or on real quantum hardware, look at the learned
        attention weights, and try continuous fingerspelling.
\end{itemize}

All code, configs, results and the paper are in the project repository
(\texttt{hybrid-qnn-sign-private}).

\end{document}
